\section{Methods}

\textbf{Targets and preprocessing.} We used RCSB PDB structures (C$\alpha$ traces) for 1UBQ(A), 2CI2(I), 5PTI(A), 1EMA(A), 1AKE(A), 1TIM(A), and a 250-residue tail segment of 2O5P chain A (residues 566–815). Unless otherwise stated, we remove near-neighbour pairs $|i-j|\le 2$ from contact/distogram scoring.

\textbf{Large-scale PDB benchmark protocol (N=1000).} To stress-test the geometric proxy claims beyond the staged suite, we constructed a reproducible RCSB sample of 1000 protein-only experimental structures (X-ray diffraction or electron microscopy; resolution $\le 2.5$~\AA), then selected the longest-chain C$\alpha$ trace per entry and filtered to lengths $60\le N\le 350$. mmCIF files are cached locally and a fixed seed controls the shuffled candidate list to ensure reproducibility (see \texttt{data/processed/pdb\_bench/pdb\_ids\_n1000\_seed0.txt} and the audit manifest in \texttt{docs/runs/full1k\_hybrid\_auric/pdb\_download\_manifest\_n1000\_seed0.csv}). All benchmark scripts live under \texttt{scripts/pdb\_bench/} and emit summaries under \texttt{data/processed/pdb\_bench/} with companion reports in \texttt{artifacts/reports/}.

\textbf{6D icosahedral embedding.} A 6D integer walk $n_i$ is projected to 3D by $x_i = s\cdot(B_{\parallel} n_i)$ followed by Kabsch alignment to the target. Perpendicular-space coordinates are $y^{\perp}_i = (B_{\perp} n_i) - w_0$. We report $\mathrm{ph}_{\max}$ as the maximum deviation from the current window proxy (and $\mathrm{ph\_ratio}=\mathrm{ph}_{\max}/\theta$). The acceptance window model is an operational proxy suitable for audit tests; future work will replace it with a full triacontahedral domain.

\textbf{Acceptance-window proxy and projector repair.} Let $W\subset\mathbb{R}^3$ denote the (proxy) acceptance window in perpendicular space. Given $y_i^\perp$, define a pointwise violation $v_i=\mathrm{dist}(y_i^\perp,W)$ and summary statistics $\mathrm{ph}_{\max}=\max_i v_i$ and $\mathrm{ph}_{\mathrm{rms}}=\sqrt{\frac{1}{N}\sum_i v_i^2}$. Under a filter-style gate, a proposal is rejected when $\mathrm{ph}_{\max}>\theta$ (or $\theta^\star$). Under a projector-style repair, we instead update the slice offset by a small inner-loop solve, e.g.\ $w_0^{\mathrm{new}}\in\arg\min_w \sum_i \mathrm{dist}(B_\perp n_i-w, W)^2$ (or a max-violation analogue), then re-evaluate feasibility under the updated $w_0$.

\textbf{Alphabets.} axis12 uses ±$e_i$. pair72 adds ±$e_i$±$e_j$ moves filtered to match the target bond length under projection. triple232 further adds ±$e_i$±$e_j$±$e_k$. In oracle-codec experiments, the step direction at residue i is assigned by nearest-neighbour quantization to the alphabet.

\textbf{Phason proxy statistics (real vs controls).} For the large-scale phason survey, we first recover a 6D walk by per-bond oracle direction quantization under triple232: for each bond direction $\Delta x_i/\|\Delta x_i\|$, we choose the alphabet direction with maximal cosine similarity under $B_{\parallel}$ and then scale that unit direction to the observed bond length $\|\Delta x_i\|$. From the induced 6D walk $n_i$, we compute $y_i^{\perp}=B_{\perp}n_i$ and set $w_0=\mathrm{mean}(y^{\perp})$; we report $\mathrm{ph}_{\mathrm{rms}}=\sqrt{\mathbb{E}\|y_i^{\perp}-w_0\|_2^2}$ and $\mathrm{ph}_{\max}=\max_i\|y_i^{\perp}-w_0\|_2$. Controls include (i) random 3D chains with the same per-bond length sequence and (ii) perturbed-native chains generated by adding small noise to bond directions while preserving bond lengths. Effect sizes are summarized via per-protein Cliff's $\delta$ and bootstrap confidence intervals (\texttt{data/processed/pdb\_bench/phason\_stats\_summary\_n1000\_seed0\_triple232.csv}).

\textbf{Multi-scale symbolic certificate metrics (Fold\_m).} To increase protocol resolution beyond a single scalar proxy, we compute multi-scale symbolic certificates from binary readouts $\rho(t)\in\{0,1\}$. Given a readout stream, we form length-$m$ sliding windows and stabilize each window via $\mathrm{Fold}_m$ into a golden-mean legal type (no adjacent 11), producing a stable-type stream. We report (i) stable-type entropy $H(\mathrm{type})$ and (ii) an SMB-style entropy-rate proxy $\widehat{h}_{\mathrm{SMB}}$ estimated from block entropies. For the full1k hybrid protocol, $\rho_A$ is defined from perpendicular-space $y^\perp$ using a \textbf{per-protein PCA axis} computed once from the real chain and reused for all null replicates, with robust median thresholding; $\rho_B$ is defined from the 6D step stream $\Delta n$ using parity/majority-sign binarization. In addition to random-chain controls, we use bond-direction shuffle and block-shuffle controls (\texttt{--shuffle-reps}, \texttt{--blockshuffle-reps}, \texttt{--blockshuffle-k}) to directly test long-range correlation structure.

\textbf{Decoy hard-negative benchmark.} We evaluated the geometry-based certificate on Decoys 'R' Us (\texttt{4state\_reduced} and \texttt{lmds}), standard decoy families containing near-native but incorrect folds. Decoys are downloaded and cached locally (not committed) under \texttt{data/raw/decoy\_cache/}, and lightweight native/decoy manifests are emitted under \texttt{docs/runs/} (including per-target histograms and summary tables). Certificate metrics are computed for each native and its decoys using a protocol-fixed axis derived \emph{exclusively} from the native $y^\perp$ and reused for all decoys of that target.

\subsubsection*{Strict axis selection protocol for decoys}
To rigorously exclude any possibility of information leakage from the decoy set into the projection axis definition, we implemented two native-only axis selection strategies:
\begin{enumerate}
    \item \textbf{Shared-PCA (Default):} the projection axis $u$ is the first principal component of the \textit{native} structure's perpendicular trace $y_{\perp}$. This single axis is frozen and applied to all decoys for that target.
    \item \textbf{Best-Native-Axis (Ablation):} as a stricter control, we select an axis $u$ from a small, pre-registered candidate family (including canonical $x/y/z$ basis vectors) by minimizing the type-entropy of the \textit{native} chain under $\rho_A$ at a fixed scale ($m=8$). This selection depends \textit{exclusively} on the native structure's internal geometry and is blind to the decoy population.
\end{enumerate}
We report results primarily on Shared-PCA, with Best-Native-Axis used to verify that separation robustness is not an artifact of axis overfitting or protocol drift.

\textbf{Blind sprint prototypes (contact/distogram + certificates).} To test whether certificates can steer search (not only discriminate post hoc), we implemented small blind-sprint prototypes under \texttt{scripts/blind\_sprint/}. These runs do not use oracle direction guidance and treat TM-score as evaluation-only, optimizing native-derived contact or sparse-distogram objectives and optionally augmenting the score with certificate penalties. Audit outputs and manifests are stored under \texttt{docs/runs/}.

\textbf{Physical consistency audit (post-hoc; CA-level).} To provide a low-cost physics-facing sanity check without modifying the geometry-driven search, we compute two coarse-grained proxies from C$\alpha$ traces: (i) a bond-length MSE relative to 3.8~\AA{} and (ii) a steric-clash proxy counting non-neighbour C$\alpha$ pairs closer than 3.5~\AA{} (reported per-residue). We combine them into a single normalized physics score: clashes/$N$ + 1000$\times$bond MSE. These metrics are intended to flag gross non-physical geometry rather than to replace full-atom energy evaluation. We compare the physics score against the geometry-based certificate ($\rho_A$ type entropy at $m=8$) and report the dual-filter funnel scatter in Fig.~S\ref{fig:physics_audit}. Scripts and lightweight CSV outputs are provided under \texttt{scripts/audit/} and \texttt{docs/runs/physics\_audit/}.

\textbf{QUARK / I-TASSER homology-ablation benchmark.} To contextualize PhasonFold's certificate-based benchmarks against widely used external structure-prediction pipelines, we prepared a small server benchmark on six short single-domain targets (all $<200$ aa; reused from our decoy narrative: \texttt{1R69}, \texttt{2CRO}, \texttt{4PTI}, \texttt{1CTF}, \texttt{1DTK}, \texttt{1SHF-A}). For each target we submit four jobs: QUARK-A (default), QUARK-B (exclude fragments from proteins with $>$30\% identity), I-TASSER-A (default), and I-TASSER-B (exclude homologous templates). We then quantify (i) cross-method topology agreement (QUARK-A model1 vs I-TASSER-A model1) and (ii) homology sensitivity within each method (A vs B). The auditable run directory (inputs, job-ID tracker, and TM-align summary CSV template) is \texttt{docs/runs/quark\_itasser\_homology\_ablation/}; large server tarballs are cached locally under \texttt{data/raw/} and not committed.

\textbf{Rosetta constrained-relax ``physical audit'' (FastRelax; PyRosetta).} To evaluate whether Omega C$\alpha$ traces occupy attractive basins of a standard all-atom energy function, we run a coordinate-constrained relaxation audit using PyRosetta (Rosetta \texttt{ref2015} score function). For each target, we first build a full-atom pose from the target sequence (\texttt{fa\_standard}) and add harmonic coordinate constraints on all C$\alpha$ atoms to the input C$\alpha$ trace (restraint strength controlled by \texttt{--coord-sd} and \texttt{--coord-weight}; main experiments use \texttt{sd=1.0\AA} and weight 1.0). We then run \textbf{FastRelax} for $nstruct$ independent trajectories (main experiments: $nstruct=20$), recording the final total Rosetta energy (REU) and the post-relax C$\alpha$ RMSD to the input trace (``drift''). We report best-of-$nstruct$ and median energy summaries and write per-trajectory CSVs used for Fig.~\ref{fig:5}B boxplots. The implementation is \texttt{scripts/benchmarks/pyrosetta\_relax\_validate.py}; lightweight outputs are stored under \texttt{docs/runs/quark\_itasser\_homology\_ablation/} (summary and per-trajectory CSVs), while large relaxed PDBs are cached under \texttt{data/raw/}.

\textbf{Chain-consistent 6D beam-lift benchmark.} To obtain a chain-consistent lift closer to operator semantics, we implement a beam search over the pair72 alphabet. At each residue, we preselect top-$k$ candidate steps by cosine similarity to the native bond direction, then run a beam over 6D states minimizing an angular mismatch term plus a small perpendicular-space regularizer $w_{\perp}\|y^{\perp}\|^2$ (for stability). The resulting 3D reconstruction is Kabsch-aligned to the native chain and scored by TM-score; we also report $\mathrm{ph}_{\mathrm{rms}}$ and compare against random-chain controls (\texttt{data/processed/pdb\_bench/beam\_lift\_bench\_summary\_n1000\_seed0\_pair72\_beam8.csv}).

\textbf{$\theta_{\mathrm{adaptive}}$ and $\theta^\star$.} $\theta_{\mathrm{adaptive}}$ is estimated from early free sampling (first 20\% of steps) as a quantile of the ph\_max distribution. $\theta^\star$ is a per-target feasible-tight threshold derived from a counterfactual feasibility floor: it is set slightly above the minimum ph\_max achievable by projector repair.

\textbf{Filter vs projector vs piecewise projector.} Filter rejects proposals with ph\_max>$\theta$. Projector performs a small inner-loop search over $w_0$ to reduce ph\_max and accepts the repaired state if feasible. Piecewise projector applies $w_0$ optimization per segment, allowing slow drift across long chains.

\textbf{Metrics.} Contact\_F1 uses an 8~\AA{} cutoff (or 12~\AA{} where stated) on C$\alpha$ distances; distogram RMSE is computed on sampled residue pairs; TM-score is computed after alignment. For the membrane segment (2O5P tail250), we compute a hydrophobic belt diagnostic by correlating residue hydropathy with radius-to-axis in a fixed axis frame estimated from the native structure.

\textbf{Reproducibility and audit.} All experiments emit audit logs containing the operator stream, accept/reject decisions, phason statistics, repair actions ($\Delta w_0$), and per-step metrics. All artifacts required to reproduce the figures and tables in this manuscript are included in the accompanying evidence repository.
